\section*{Chapter 7: The First Image Of A Black Hole}
\begin{description}[style=multiline,leftmargin=2em,labelwidth=1.5em]

% DOCX paragraph 2071
\item[1.]\hypertarget{note:07-01}{}Plasma\index{plasma|textbf} is one of the four states of matter, together with solid,
liquid, and gaseous. Matter in this state is characterised by a mixture of
ions\index{ions} (i.e., atoms stripped of some or all of their electrons) and
free electrons. The plasma state is widespread in astrophysics due to the very high
energies that can be reached in the cosmos, for example on a star's surface.

% DOCX paragraph 2072
\item[2.]\hypertarget{note:07-02}{}H. Falcke, M. Kramer, L. Rezzolla, ERC Synergy Grant \emph{BlackHoleCam: Imaging the
Event Horizon of Black Holes}; Grant No. 610058 (2013-2021).

% DOCX paragraph 2073
\item[3.]\hypertarget{note:07-03}{}The impact parameter is used in particle physics and gravitational physics, to measure
the perpendicular distance between a target and the closest point reached by a particle
along its trajectory.

% DOCX paragraph 2074
\item[4.]\hypertarget{note:07-04}{}Obviously, given the spherical symmetry\index{spherical symmetry} of a Schwarzschild black
hole\index{black hole!Schwarzchild}, the unstable circular photon orbit is an entire spherical
surface with a radius equal to the radius of the light ring. The concept of the light
ring can also be extended to rotating black holes, but in this case, the picture becomes
much more complicated. In addition to the distinction between photon trajectories that
are \emph{corotating} and \emph{counter-rotating} with respect of the black hole, the
spherical surface of the light ring needs to be replaced by a region that is only
approximately spherical.

% DOCX paragraph 2075
\item[5.]\hypertarget{note:07-05}{}R. Penrose and R. M. Floyd, in \emph{Nature Physical Science}, 229, 177 (1971).

% DOCX paragraph 2076
\item[6.]\hypertarget{note:07-06}{}The effect just described takes the name of \emph{longitudinal} \emph{(relativistic)
Doppler effect}\index{relativistic!Doppler effect}, to distinguish it from the \emph{transverse
relativistic Doppler effect}, which leads to a variation of the frequency even when the
motion is in a direction perpendicular to that of the emission of light. However, the
transverse Doppler effect does not play a major role in what we are discussing here.

% DOCX paragraph 2077
\item[7.]\hypertarget{note:07-07}{}In mathematics, all these properties are expressed through the \emph{degree of
nonlinearity} of the equations, which indicates how a small variation in one of the
properties of the plasma can lead to enormous variations in its behaviour.

% DOCX paragraph 2078
\item[8.]\hypertarget{note:07-08}{}A positron is the antimatter counterpart of an electron and is, therefore, a particle
with the same rest mass of an electron but with a positive charge.

% DOCX paragraph 2079
\item[9.]\hypertarget{note:07-09}{}There are a number of scenarios that could potentially provide the energy for the
acceleration. One possibility is that it comes directly from the black hole's rotational
energy, through what is known as the \emph{Blandford-Znajek process} (see R. D.
Blandford and R. L. Znajek, in \emph{Monthly Notices of the Royal Astronomical Society},
179, 433 (1977)). Another possibility is that it is instead supplied by the rotational
energy of the disc, through what is known as the \emph{Blandford-Payne process} (see R.
D. Blandford and D. G. Payne, in \emph{Monthly Notices of the Royal Astronomical
Society}, 199, 883 (1982)). Finally, there is a scenario in which the acceleration is
provided through dissipative processes within the magnetised plasma.

% DOCX paragraph 2080
\item[10.]\hypertarget{note:07-10}{}T. Karras, S. Laine, T. Aila, \emph{A Style-Based Generator Architecture for Generative
Adversarial Networks}, arXiv: 1812.04948.

% DOCX paragraph 2081
\item[11.]\hypertarget{note:07-11}{}Y. Mizuno, Z. Younsi, CM Fromm, O. Porth, M. De Laurentis, H. Olivares, H. Falcke, M.
Kramer, L. Rezzolla, in ``Nature Astronomy'', 2, 585 (2018); H. Olivares, Z. Younsi, C.
M. Fromm, M. De Laurentis, O. Porth, Y. Mizuno, H. Falcke, M. Kramer, L. Rezzolla, in
\emph{Monthly Notices of the Royal Astronomical Society}, 497, 521 (2020)

\end{description}
